In this thesis, we have explored various use cases of \glspl{llm} in the context
of local \gls{cli} assistants, with constraings on computational resources.
The first use case was the usage of \glspl{llm} in a role of command versus
natural language classificator, to make the integration of the assistant into
the system shell seamless, but this approach was not successful. Additionally,
there are much better approaches to this problem, that use other \glspl{ml}
models, that are much more efficient and accurate for this task.
The second use case was the usage of embedding \glspl{llm} for retrieval of
relevant information from a database, to make the search for relevant in local
documents (in scope of this thesis --- man manual pages) more efficient and
accurate. We have explored three approaches to this problem: Document-level,
Passage-level, and Document-level with Passage-level reranking. Overall, for the
specific use case ot this theis, the best approach is Document-level with
Passage-level reranking for biger document collections, and passage-level for
smaller ones, as it provides the best balance between latency and accuracy. The
use of Document-level approach is not recommended, as it generates much more
output, and it is much more difficult to control the size of the generated
output, as it depends on the size of the documents.
The third use case was the usage of \glspl{llm} for generation of natural
language responses to user queries, based on the retrieved information from the
previous use case. This approach is very effective, as it allows to generate
more natural and concise responses in short time, but it showed to be very
sensitive to the quality and size of the retrieved information, which affected
the choice of the best approach for the previous use case.
